\documentclass[10pt,a4paper]{amsart}
\usepackage{amsmath,amssymb,amsthm,graphicx}
\usepackage{verbatim}
\newenvironment{reference}{\comment}{\endcomment}
%\newenvironment{reference}{}{}
\newenvironment{slogan}{\comment}{\endcomment}
\newenvironment{history}{\comment}{\endcomment}

% For commutative diagrams we use Xy-pic
\usepackage[all]{xy}

% We use 2cell for 2-commutative diagrams.
\xyoption{2cell}
% Theorem environments.
%
\theoremstyle{plain}
\newtheorem{theorem}[subsection]{Theorem}
\newtheorem{proposition}[subsection]{Proposition}
\newtheorem{lemma}[subsection]{Lemma}

\theoremstyle{definition}
\newtheorem{definition}[subsection]{Definition}
\newtheorem{example}[subsection]{Example}
\newtheorem{exercise}[subsection]{Exercise}
\newtheorem{situation}[subsection]{Situation}

\theoremstyle{remark}
\newtheorem{remark}[subsection]{Remark}
\newtheorem{remarks}[subsection]{Remarks}

\numberwithin{equation}{subsection}

% Macros
%
\def\lim{\mathop{\mathrm{lim}}\nolimits}
\def\colim{\mathop{\mathrm{colim}}\nolimits}
\def\Spec{\mathop{\mathrm{Spec}}}
\def\Hom{\mathop{\mathrm{Hom}}\nolimits}
\def\Ext{\mathop{\mathrm{Ext}}\nolimits}
\def\SheafHom{\mathop{\mathcal{H}\!\mathit{om}}\nolimits}
\def\SheafExt{\mathop{\mathcal{E}\!\mathit{xt}}\nolimits}
\def\Sch{\mathit{Sch}}
\def\Mor{\mathop{\mathrm{Mor}}\nolimits}
\def\Ob{\mathop{\mathrm{Ob}}\nolimits}
\def\Sh{\mathop{\mathit{Sh}}\nolimits}
\def\NL{\mathop{N\!L}\nolimits}
\def\CH{\mathop{\mathrm{CH}}\nolimits}
\def\proetale{{pro\text{-}\acute{e}tale}}
\def\etale{{\acute{e}tale}}
\def\QCoh{\mathit{QCoh}}
\def\Ker{\mathop{\mathrm{Ker}}}
\def\Im{\mathop{\mathrm{Im}}}
\def\Coker{\mathop{\mathrm{Coker}}}
\def\Coim{\mathop{\mathrm{Coim}}}

% Boxtimes
%
\DeclareMathSymbol{\boxtimes}{\mathbin}{AMSa}{"02}

%
% Macros for moduli stacks/spaces
%
\def\QCohstack{\mathcal{QC}\!\mathit{oh}}
\def\Cohstack{\mathcal{C}\!\mathit{oh}}
\def\Spacesstack{\mathcal{S}\!\mathit{paces}}
\def\Quotfunctor{\mathrm{Quot}}
\def\Hilbfunctor{\mathrm{Hilb}}
\def\Curvesstack{\mathcal{C}\!\mathit{urves}}
\def\Polarizedstack{\mathcal{P}\!\mathit{olarized}}
\def\Complexesstack{\mathcal{C}\!\mathit{omplexes}}
% \Pic is the operator that assigns to X its picard group, usage \Pic(X)
% \Picardstack_{X/B} denotes the Picard stack of X over B
% \Picardfunctor_{X/B} denotes the Picard functor of X over B
\def\Pic{\mathop{\mathrm{Pic}}\nolimits}
\def\Picardstack{\mathcal{P}\!\mathit{ic}}
\def\Picardfunctor{\mathrm{Pic}}
\def\Deformationcategory{\mathcal{D}\!\mathit{ef}}
\usepackage{cleveref}
\crefname{originaltheorem}{original theorem}{original theorems}
\Crefname{originaltheorem}{Original theorem}{Original theorems}
\makeatletter
\newlabel{theorem-K-injective-embedding-grothendieck}{{079P}{source}{}{sourcepointer-1}{}}
\newlabel{theorem-K-injective-embedding-grothendieck@cref}{{[originaltheorem][0][]079P}{[1][0][]source}{}{}{}}
\makeatother
\crefname{originallemma}{original lemma}{original lemmas}
\Crefname{originallemma}{Original lemma}{Original lemmas}
\makeatletter
\newlabel{derived-lemma-product-K-injective}{{0BK6}{source}{}{sourcepointer-2}{}}
\newlabel{derived-lemma-product-K-injective@cref}{{[originallemma][0][]0BK6}{[1][0][]source}{}{}{}}
\makeatother
\title{Direct sums and products in the derived category of a Grothendieck abelian category}
\author{The Stacks Project Authors}
\date{stacks-project:injectives.tex:local-snapshot-20261012; GNU FDL 1.2 or later}
\begin{document}
\maketitle
\noindent\textit{Editorial note.} All delivered mathematical text is the author's own, copied from the pinned Stacks Project chapter file sites-cohomology.tex; no statement or proof step was rewritten, reordered or removed.\par\medskip
\noindent\textit{Editorial note.} The proof is the author's complete argument for both assertions, the direct sum case being carried out in the span and the product case being reduced to a named derived-category result carried on this card as a navigation stub.\par\medskip
\noindent\textit{Editorial note.} Every labelled reference occurring inside the counted spans is carried in assembly.external\_references as a non-load-bearing navigation pointer (same-chapter ones at their definition site, cross-chapter ones at the invocation inside the counted span), so no load-bearing mathematics is replaced by a stub. No bibliography citation occurs inside the counted spans.\par\medskip
\noindent\textit{Editorial note.} The same-name check for this card scanned every staged Stacks chapter of this repository -- all sources/stacks-*/pinned-original-source.tex files, the card's own chapter excluded -- and the labels found are recorded in dedup.same\_name\_cross\_chapter on the card.\par\medskip
\noindent\textit{Editorial note.} The preamble is the chapter's own retained minimal preamble (sources/stacks-sites-cohomology/retained-preamble.tex), a verbatim excerpt of the native preamble: the Xy-pic setup and the native theorem declarations with the mathematical macros.\par\medskip
\noindent\textit{Allowed standard prerequisites (editorial).}\begin{itemize}
\item Products of K-injective complexes. Derived Categories, Lemma derived-lemma-product-K-injective. Supplies that the termwise product of K-injective complexes is a product in the derived category, which is the second assertion.
\item K-injective resolutions. Same chapter, Theorem theorem-K-injective-embedding-grothendieck and the discussion of K-injective resolutions. Supplies that every object of the derived category can be represented by a K-injective complex, which is what makes the construction available.
\item Termwise direct sums and colimits. Derived Categories, the description of direct sums in the derived category as termwise direct sums. Supplies the first assertion about direct sums of complexes.
\end{itemize}\medskip
\ifdefined\hypertarget\hypertarget{sourcepointer-1}{}\fi
\noindent\textit{Source reference.} Original theorem 079P: sources/stacks-injectives/pinned-original-source.tex, source lines 2024--2024. Navigation pointer invoked at proof line 2212; the counted span contains the pointer only and the cited result is not reproduced inside it, so the pointer is registered as a non-load-bearing navigation stub in the established form for Stacks material. No bibliography citation is involved.\par\medskip
\ifdefined\hypertarget\hypertarget{sourcepointer-2}{}\fi
\noindent\textit{Source reference.} Original lemma 0BK6: sources/stacks-derived/pinned-original-source.tex, source lines 9917--9917. Navigation pointer invoked at proof line 2259; the counted span contains the pointer only and the cited result is not reproduced inside it, so the pointer is registered as a non-load-bearing navigation stub in the established form for Stacks material. No bibliography citation is involved.\par\medskip
\begin{lemma}
\label{lemma-derived-products}
Let $\mathcal{A}$ be a Grothendieck abelian category.
Then
\begin{enumerate}
\item $D(\mathcal{A})$ has both direct sums and products,
\item direct sums are obtained by taking termwise direct sums of
any complexes,
\item products are obtained by taking termwise products of
K-injective complexes.
\end{enumerate}
\end{lemma}

\begin{proof}
Let $K^\bullet_i$, $i \in I$ be a family of objects of $D(\mathcal{A})$
indexed by a set $I$. We claim that the termwise direct sum
$\bigoplus_{i \in I} K^\bullet_i$ is a direct sum in $D(\mathcal{A})$.
Namely, let $I^\bullet$ be a K-injective complex. Then we have
\begin{align*}
\Hom_{D(\mathcal{A})}(\bigoplus\nolimits_{i \in I} K^\bullet_i, I^\bullet)
& =
\Hom_{K(\mathcal{A})}(\bigoplus\nolimits_{i \in I} K^\bullet_i, I^\bullet) \\
& =
\prod\nolimits_{i \in I} \Hom_{K(\mathcal{A})}(K^\bullet_i, I^\bullet) \\
& =
\prod\nolimits_{i \in I} \Hom_{D(\mathcal{A})}(K^\bullet_i, I^\bullet)
\end{align*}
as desired. This is sufficient since any complex can be represented
by a K-injective complex by
Theorem \ref{theorem-K-injective-embedding-grothendieck}.
To construct the product, choose a K-injective resolution
$K_i^\bullet \to I_i^\bullet$ for each $i$. Then we claim that
$\prod_{i \in I} I_i^\bullet$ is a product in $D(\mathcal{A})$.
This follows from
Derived Categories, Lemma \ref{derived-lemma-product-K-injective}.
\end{proof}

\end{document}
